\chapter{Introducción}

El estudio de las conexiones en transformadores trifásicos es un pilar fundamental en la formación sobre sistemas
eléctricos de potencia. Las diferentes configuraciones de conexión en los bobinados primario y secundario (tales como
estrella, triángulo y zig-zag) introducen desfases angulares específicos entre las tensiones de línea de ambos lados.
Dicho desfase se representa mediante el denominado índice horario, un valor numérico entero del 0 al 11 que, multiplicado
por $30^\circ$, indica el retraso angular de los vectores del secundario respecto a los del primario.

Para comprender y verificar este fenómeno de manera didáctica y visual, en este trabajo práctico se utiliza la plataforma
virtual y educativa Aula Moisan de la Universidad de Valladolid. A través de su simulador interactivo, se analizan y
determinan las configuraciones de índice horario y sus diagramas vectoriales correspondientes.

\section{Objetivos}
\begin{itemize}
    \item Determinar de forma gráfica y experimental, mediante el simulador de Aula Moisan, el índice horario resultante
    de diversas configuraciones de conexión en transformadores trifásicos.
    \item Estudiar las características de las conexiones en estrella, triángulo y zig-zag en los devanados de alta
    tensión (AT) y baja tensión (BT).
    \item Analizar la influencia de la polaridad de las conexiones y el orden de sucesión de fases en el desfasaje angular.
    \item Comprender la importancia práctica del índice horario, especialmente para el acoplamiento en paralelo de
    transformadores y la mitigación de problemas como armónicos y desequilibrios de carga.
\end{itemize}
